\section{Weekly Summary}

\begin{tcolorbox}[colback=yellow!10,colframe=black!50,title=AI-Generated Content]
\noindent The summary below was written by Claude (Anthropic) from that week's regression outputs and series descriptions. It has not been reviewed or fact-checked by a human, and should be read as an automated, informal recap -- not analysis.
\end{tcolorbox}

Welcome back to the corner of the internet where we throw two unrelated economic time series into a blender, hit "regress," and pretend we're doing science. This week was a modest two-day affair, which is to say we generated exactly two chances to embarrass ourselves and, frankly, delivered on both.

Monday, August 18, gave us the crown jewel of trend-mirages: a discontinued index of hospital and nursing-home output pitted against house prices in Nassau and Suffolk County, New York. The raw regression came out swinging -- an R-squared of 0.60, a p-value with more zeros after the decimal than I care to count, the whole nine yards. For one glorious moment I nearly believed that the fate of Long Island real estate is written in the ledgers of nursing homes. And then we detrended it, and the entire relationship evaporated like a puddle in July. R-squared cratered to 0.08, the p-value ballooned to a limp 0.23, and the "significant" link revealed itself for what it always was: two things that both went up over the years because that's what a lot of things do. Textbook drift. Chef's kiss. This is the exact species of correlation this project exists to catch and gently mock.

Tuesday, August 19, mercifully spared us the whole rise-and-fall drama by simply not delivering much of anything. U.S. Government Grants Excluding Military versus a mining output price deflator managed a raw p-value of 0.08 -- not even significant by the usual standards, just close enough to make you squint. Once detrended, it slumped to a p-value of 0.53 and an R-squared of 0.016, which is statistician-speak for "these two numbers have never met and would not recognize each other on the street." A blessedly boring result, and honestly a relief after Monday's emotional rollercoaster.

The obligatory reminder, because my lawyer insists: these series are paired at random, nobody is claiming that hospital output moves house prices or that government grants haunt the mining sector, and the entire point of the exercise is that spurious correlation is the default setting of the universe, not some rare anomaly. This week the detrending did its job with flawless consistency by demolishing everything, which means we found zero genuine relationships and had a lovely time doing it. See you next week for more of the same.
